\subsection{Fase 4: Análise e Julgamento da Avaliação}
\label{sec:fase4-analise}

\subsubsection{Execução Documentada e Elegibilidade}

O corte documental distingue resultados oficiais de relatos preliminares. Na \texttt{main} (commit \texttt{da717f2}), os CSV brutos estavam vazios e o cálculo gerava 25 registros inconclusivos e um não aplicável (M3.1.2 no aplicativo). O piloto S01, de 20/11, não integra o cálculo: faltam aquecimento, marcações independentes e identificação do instrumento de volume.

A branch \texttt{feat/pedro-s03-fase4} (commit \texttt{14ab804}) acrescenta a sessão S03 de 21/11, resumida na Tabela~\ref{tab:fase4-observacoes}; a conferência independente das 14 linhas está pendente. Os termos de consentimento foram aplicados, com vias privadas; a conciliação dos códigos de veículos segue pendente (fontes em \texttt{docs/fase4-analise.md}).

\begin{table}[htbp]
\centering
\caption{Observações preliminares de S03/P6; não são níveis oficiais.}
\label{tab:fase4-observacoes}
\footnotesize
\begin{tabular}{P{2.1cm}P{3.2cm}P{7.2cm}}
\toprule
Contexto & Tentativas elegíveis relatadas & Resultado e limite \\
\midrule
CT01: portal/Chrome & 3 conclusões em 3 & Cinco erros de instrumento preservados à parte; execução automatizada. \\
CT03: portal/Android & 2 conclusões em 3 & Uma sessão expirada durante a consulta, preservada como falha. \\
CT05: aplicativo & 0 conclusões em 3 & Erro ao buscar veículos; Aplicativo 2.6.3 e aparelho superior ao planejado (D006). \\
\bottomrule
\end{tabular}
\par\smallskip{\footnotesize Fonte: registro S03, commit \texttt{14ab804}; conferência independente pendente.}
\end{table}

CT02, CT04 e CT06 não têm execução nesse relato. O impedimento de login no Firefox ocorreu fora das tentativas e não foi contado como falha do portal. As durações dos scripts incluem esperas fixas e não substituem P4; P6 não mede as taxas por rede de P7. Versões e aparelhos diferentes, automação e ausência de pares temporais válidos impedem inferir consistência, causa técnica ou superioridade geral.

\subsubsection{Continuação com o Caso V4 (S06 e S08)}
\label{subsec:fase4-v4}

Em 22 e 23/11/2026, a coleta prosseguiu com outro titular consentido (caso V4), no computador D06 e no Samsung SM-S731B (Aplicativo 2.7.0), com início em \emph{Consultar}. A Tabela~\ref{tab:fase4-v4} traz a primeira leitura automática pelos quadros, sem aquecimentos; sem a segunda leitura independente, nenhum nível é atribuído.

\begin{table}[!htbp]
\centering
\caption{Observações provisórias de S06/S08 (caso V4): mediana e P90; não são níveis oficiais.}
\label{tab:fase4-v4}
\scriptsize
\begin{tabular}{P{3.0cm}P{4.6cm}P{4.6cm}}
\toprule
Métrica e rede & Portal & Aplicativo \\
\midrule
M2.1.1, Wi-Fi & 10/10; 0,145 s; 0,683 s & 10/10; 0,774 s; 0,858 s \\
M2.1.2, Wi-Fi & 10/10; 0,136 s; 0,153 s & 10/10; 0,641 s; 0,816 s \\
M2.2.1, Wi-Fi & 0,0107 MB (navegador) & 0,0649 MB (UID) \\
M3.2.3, Wi-Fi & 10/10; 0,083 s; 0,083 s & 10/10; 0,408 s; 0,667 s \\
M3.2.3, 4G & 10/10; 0,142 s; 0,183 s & 10/10; 0,726 s; 0,942 s \\
M3.2.3, degradada & 8/10; 0,850 s; 1,500 s & Não avaliada (D021) \\
\bottomrule
\end{tabular}
\par\smallskip{\footnotesize Fonte: \texttt{dados/brutos/} e registros S06/S08; conferência independente pendente.}
\end{table}

Os escopos de volume diferem e não sustentam superioridade. No Aplicativo, a lista fica abaixo da área visível; o B04 foi excluído por esperas do instrumento, e o B09 usou rolagens rápidas do aparelho desde o toque, cujo curso integra o tempo (D026). Falhas e tempo limite da rede degradada permanecem no denominador. Essa condição não pôde ser imposta ao Aplicativo: o adaptador não sustentou ponto de acesso, e o Aplicativo ignorou proxy e recusou VPN reversa (D021).

\subsubsection{Respostas GQM e Relações entre Características}

Q1.1, Q1.2 e Q1.3 permanecem inconclusivas: não há registros oficiais de completude, consistência e fluxo. Q2.1 e Q2.2 têm amostras completas de V4 (Tabela~\ref{tab:fase4-v4}), mas sem a segunda leitura exigida não recebem nível. Q3.1 tem um relato de bloqueio no aplicativo, sem conferência; Q3.2 tem tentativas do portal nas três redes e do aplicativo em Wi-Fi e 4G, também pendentes. Assim, as sete hipóteses não recebem julgamento definitivo e G1, G2 e G3 permanecem inconclusivos nos dois canais. A ausência de dados ou de conferência não equivale a paridade nem a aprovação.

A falha antes da lista em CT05 impede observar campos e detalhes, relacionando cobertura de contexto e adequação funcional; não se atribui completude zero a uma tela não alcançada. A expiração em CT03 mostra por que falhas devem acompanhar tempos, evitando considerar apenas consultas concluídas. As primeiras leituras de V4 não sustentam ainda associações entre volume, tempo e completude; na rede degradada, só o portal foi observado. Cada questão exige todas as métricas aplicáveis no mínimo aceitável e cada objetivo exige todas as suas questões atendidas.

\subsubsection{Ação de Melhoria Entregue}
\label{subsec:acao-melhoria}

O bloqueio relatado em CT05 motivou uma proposta de processo com templates técnicos, tipo aceito na issue \#36. O artefato em \texttt{melhorias/consulta-infracoes/} entrega ficha de incidente, matriz de quatro cenários e inspeção automática: conferir a evidência, reproduzir a configuração, separar intervenções por versão, preparar diagnóstico e repetir três tentativas P6 por contexto. Apoia D4 (suporte por plataforma) e D1 (investimento), priorizando conclusão da tarefa sem avaliar usabilidade.

A aprovação do contexto preserva o mínimo de duas conclusões em três; três em três é apenas meta local de regressão. A proposta foi incorporada ao repositório, mas não altera o serviço do Detran-DF. Revisão humana e eficácia sobre o produto permanecem pendentes. Não se declara defeito corrigido, causa identificada ou encaminhamento ao órgão.
